\documentclass[11pt]{article}

\usepackage[a4paper,margin=1in]{geometry}
\usepackage{amsmath,amssymb}
\usepackage{graphicx}
\usepackage{hyperref}
\usepackage{url}

\title{GRA Swarm Arena: An Abstract Multi-Agent Simulation for Entropy-Guided Coordination}
\author{AAA\\
\small Independent Researcher\\
\small \texttt{youremail@example.com}}
\date{2026}

\begin{document}
\maketitle

\begin{abstract}
We introduce \emph{GRA Swarm Arena}, an abstract grid-based multi-agent simulation
designed to study entropy-guided coordination in swarms of simple agents.
Agents move on a 2D arena with hazards and goal cells and may employ different
strategies, including a simplified GRA-like coordinator that nudges them
towards lower-entropy trajectories.
The environment is intentionally non-domain-specific and is intended for
educational and research use rather than any real-world military applications.
\end{abstract}

\section{Introduction}

Multi-agent systems and swarm coordination are widely studied in robotics,
control theory, and artificial intelligence.
In many scenarios, agents must navigate environments with uncertainty, hazards,
and limited information.
We propose \emph{GRA Swarm Arena} as a minimal, game-like environment to explore
how a simple ``GRA-like'' mechanism can reduce the entropy of agent trajectories
and thus improve coherence of swarm behaviour.

The project is implemented as an open-source Python repository\footnote{URL will be added upon publication.}
and is meant as a playground for experimenting with entropy-based metrics,
strategy comparison, and basic coordination schemes.

\section{Environment}

The arena is a discrete $N \times N$ grid.
Each cell can be:
\begin{itemize}
  \item empty (neutral),
  \item a \emph{trap} with an associated penalty,
  \item a \emph{goal} cell.
\end{itemize}

A configuration file (YAML) specifies:
\begin{itemize}
  \item grid size $N$,
  \item number and types of traps,
  \item start positions for agents,
  \item goal positions.
\end{itemize}

Agents start from one of the designated start cells and attempt to reach any goal cell within
a fixed number of steps.

\section{Agents and Strategies}

Each agent occupies a grid cell and can move in four directions: up, down, left, right,
subject to arena bounds.
We currently implement three basic strategies:
\begin{itemize}
  \item \textbf{Random}: choose a move uniformly at random.
  \item \textbf{Greedy}: choose a move that reduces the Manhattan distance to a goal cell.
  \item \textbf{GRA-guided}: use a hint function provided by a GRA-like coordinator.
\end{itemize}

Agents record their paths as sequences of visited cells, which allows us to compute
trajectory statistics such as path entropy.

\section{Entropy Metric}

Given a path $P = (c_1, c_2, \dots, c_T)$, where each $c_t$ is a cell on the grid,
we consider the empirical distribution over visited cells and define a simple
entropy metric
\[
H(P) = - \sum_{c} p(c) \log p(c),
\]
where $p(c)$ is the fraction of visits to cell $c$ along the path.
This quantity serves as a proxy for ``chaotic'' behaviour: high entropy corresponds
to widely spread, repetitive, or unfocused trajectories.

\section{GRA-like Coordination}

The GRA coordinator operates at game level.
At each time step, it may:
\begin{itemize}
  \item evaluate current paths of agents via the entropy metric,
  \item provide a hint function to GRA-guided agents that discourages revisiting
        heavily overused cells.
\end{itemize}

In the current implementation, the hint function selects moves that minimize
the count of visits to candidate next cells along the agent's own path.
This can be seen as a very simple form of entropy regularization on trajectories.

The GRA parameters (e.g.\ entropy weight, reset threshold) are exposed in configuration
and can be swept in experiments.

\section{Experiments}

The repository includes Jupyter notebooks for basic experiments:
\begin{itemize}
  \item \emph{Entropy vs trap density}: vary the number of traps and observe how
        mean path entropy changes.
  \item \emph{Strategy comparison}: compare random, greedy, and GRA-guided strategies
        in terms of mean path entropy and average steps to reach a goal.
  \item \emph{GRA parameter sweep}: explore how different entropy-related parameters
        affect swarm behaviour.
\end{itemize}

These experiments are intended as starting points; the environment can be extended
with more complex strategies, reward structures, or observation models.

\section{Discussion and Limitations}

GRA Swarm Arena is intentionally minimal and abstract.
It does not model real-world physics, sensors, or communications, and is not meant
to approximate any specific operational scenario.
Its main purpose is to provide a clean, reproducible sandbox where entropy-based
coordination ideas can be tested and visualized.

\section{Conclusion}

We presented GRA Swarm Arena, a lightweight open-source environment for studying
entropy-guided coordination in agent swarms.
Future work may include richer agent models, learning-based strategies, and more
sophisticated GRA-like coordinators that integrate multiple objectives.

\section*{Acknowledgements}

We thank the broader open-source community for tools and inspiration that made
this project possible.

\bibliographystyle{plain}
\bibliography{references}

\end{document}
